% ---------------------------------------------------------------------
%  PHỤ LỤC. HIỆN THỰC HÓA USE CASE (chi tiết của mục Hiện thực hóa các use
%  case trọng tâm trong Chương 3)
% ---------------------------------------------------------------------
\section{Hiện thực hóa use case}\label{pl:hien-thuc-uc}

Phụ lục này là phần chi tiết của mục~\ref{ssec:hien-thuc-uc}, theo cùng khuôn hiện thực đã nêu ở đó: mỗi use case do một dịch vụ ứng dụng của module sở hữu đảm nhận, chạy trong một giao dịch theo thứ tự khóa toàn cục và phát sự kiện miền cho các tác dụng phụ.

Bảng~\ref{tab:c3-uc-thiet-ke} là ma trận truy vết 53 use case, gom theo chín gói như Bảng~\ref{tab:danh-sach-uc}. Cột Module ghi module sở hữu dịch vụ, trong ngoặc là module tham gia cùng giao dịch hoặc được đọc qua cổng. Cột API dùng mã của mục~\ref{sec:api}, trong ngoặc là API hỗ trợ cùng luồng. Tên bảng được ghi không kèm schema vì là duy nhất trong toàn CSDL; schema của bảng suy ra từ module sở hữu nêu ở cột Module. Ở cột Sự kiện, mũi tên chỉ bên nhận, tên trong ngoặc là hàng đợi River, còn “đồng bộ” nghĩa là ghi \texttt{sync.space\_changes} và phát \texttt{space:\{sid\}} sau commit. Cột Sơ đồ ghi số hình biểu đồ tuần tự hoặc máy trạng thái trong chương này. Chỉ năm API không gắn với một use case cụ thể: A05 làm mới phiên cho mọi lời gọi, A16 xóa tài khoản theo \nfr{NFR18}, Y01 và Y02 phục vụ đồng bộ và thời gian thực (\nfr{NFR10}, \nfr{NFR11}), G01 liệt kê các không gian để ứng dụng chuyển không gian làm việc.

{\footnotesize
\begin{longtable}{|C{1.0cm}|L{2.1cm}|L{1.9cm}|L{4.0cm}|L{3.4cm}|C{\dimexpr\textwidth-12\tabcolsep-7\arrayrulewidth-1.0cm-2.1cm-1.9cm-4.0cm-3.4cm\relax}|}
\caption{Ma trận truy vết từ use case tới thiết kế backend}\label{tab:c3-uc-thiet-ke}\\
\hline
\hd{Mã UC} & \hd{Module} & \hd{API} & \hd{Bảng dữ liệu chính} & \hd{Sự kiện, tác vụ nền} & \hd{Sơ đồ}\\ \hline
\endfirsthead
\multicolumn{6}{l}{\footnotesize\itshape Bảng \thetable\ (tiếp theo)}\\ \hline
\hd{Mã UC} & \hd{Module} & \hd{API} & \hd{Bảng dữ liệu chính} & \hd{Sự kiện, tác vụ nền} & \hd{Sơ đồ}\\ \hline
\endhead
\multicolumn{6}{|l|}{\hs{Gói 1. Xác thực và quản lý tài khoản}}\\ \hline
UC01 & identity \newline (+~household) & A01, A02 (A03) & \path{users}, \path{otp_challenges}; \path{spaces}, \path{space_members} & \path{user.registered} $\to$ thư OTP (\path{notify}) & \ref{fig:seq-dang-ky-otp}, \ref{fig:trang-thai-tai-khoan}\\ \hline
UC02 & identity & A02, A03 & \path{otp_challenges}; Redis \path{otp:cooldown}, \path{otp:hourly} & gửi lại thư OTP (\path{notify}) & \ref{fig:seq-dang-ky-otp}\\ \hline
UC03 & identity & A04 (A14) & \path{users}, \path{refresh_sessions}, \path{devices} & — & \ref{fig:seq-dang-nhap-lam-moi}, \ref{fig:seq-dung-lai-ma-lam-moi}\\ \hline
UC04 & identity & A07, A02, A08 & \path{otp_challenges}, \path{users}, \path{refresh_sessions} & \path{password.reset_requested} $\to$ thư OTP & \ref{fig:seq-dang-ky-otp}\\ \hline
UC05 & identity & A06, A15 & \path{refresh_sessions}, \path{devices}; Redis \path{auth:ver} & — & —\\ \hline
UC06 & identity & A09, A10 (Y03, Y04) & \path{users}, \path{media.objects}; \path{space_changes} & \path{media.uploaded} $\to$ xử lý ảnh (\path{media}); đồng bộ mọi không gian & —\\ \hline
UC07 & identity & A11 & \path{users}, \path{refresh_sessions}; Redis \path{auth:ver} & — & —\\ \hline
UC08 & identity & A12, A13 & \path{user_settings}; Redis \path{cfg:all}, \path{suggest:{sid}} & đổi $N$ $\to$ xóa \path{suggest:{sid}} & —\\ \hline
\multicolumn{6}{|l|}{\hs{Gói 2. Quản lý nhóm gia đình}}\\ \hline
UC09 & household & G02 (G04) & \path{spaces}, \path{space_members} & \path{CopySpaceData} (\path{default}); đồng bộ & —\\ \hline
UC10 & household & G05, G06 & \path{invite_codes}, \path{spaces}; Redis \path{cfg:all} & — & —\\ \hline
UC11 & household & G07 & \path{invite_codes}, \path{spaces}, \path{space_members}; \path{users} & \path{member.joined} $\to$ trưởng nhóm (\path{notify}); đồng bộ & \ref{fig:seq-tham-gia-nhom}\\ \hline
UC12 & household & G03 & \path{spaces}, \path{space_members}; \path{users}, \path{space_changes} & — & —\\ \hline
UC13 & household \newline (+~shopping) & G08 & \path{space_members}; \path{shopping_items} & \path{member.removed} $\to$ thông báo; đồng bộ & \ref{fig:seq-roi-nhom-chuyen-quyen}\\ \hline
UC14 & household & G09 & \path{space_members}, \path{invite_codes} & \path{leader.transferred} $\to$ thông báo; đồng bộ & \ref{fig:seq-roi-nhom-chuyen-quyen}\\ \hline
UC15 & household \newline (+~shopping) & G10 & \path{space_members}, \path{spaces}; \path{shopping_items} & \path{member.left} $\to$ thông báo; \path{PurgeSpace} khi giải tán; đồng bộ & \ref{fig:seq-roi-nhom-chuyen-quyen}\\ \hline
\multicolumn{6}{|l|}{\hs{Gói 3. Quản lý danh sách mua sắm}}\\ \hline
UC16 & shopping & S02, S04 (S01, S03, S05) & \path{shopping_lists}, \path{shopping_items} & đồng bộ & —\\ \hline
UC17 & shopping & S07, S08, S09 & \path{shopping_items}, \path{shopping_lists} & đồng bộ & \ref{fig:trang-thai-mat-hang-be}\\ \hline
UC18 & shopping & S06 & \path{shopping_lists}, \path{shopping_items} & đồng bộ & —\\ \hline
UC19 & shopping & S10, S11 & \path{shopping_items}; \path{space_members} & \path{AssignmentDigest} $\to$ \path{item.assigned}, \path{item.unassigned} $\to$ FCM (\path{notify}); đồng bộ & \ref{fig:seq-phan-cong}\\ \hline
UC20 & shopping & S12, S13 & \path{shopping_items}, \path{shopping_lists} & \path{items.purchased} $\to$ đồng bộ, báo cáo & \ref{fig:seq-outbox-mua-hang}\\ \hline
UC21 & shopping \newline (+~pantry) & S14 & \path{shopping_items}; \path{food_lots}, \path{stock_movements} & \path{items.stocked}, \path{lot.changed} $\to$ đồng bộ & \ref{fig:seq-nhap-tu}\\ \hline
UC22 & mealplan \newline (+~recipe, pantry, shopping) & S15 (S01) & \path{meal_entries}; \path{recipe_ingredients}, \path{food_lots}, \path{shopping_lists}, \path{shopping_items} & đồng bộ khi tạo danh sách & \ref{fig:seq-sinh-ds-ke-hoach}\\ \hline
UC23 & shopping & S01, S03 & \path{shopping_lists}, \path{shopping_items} & — & —\\ \hline
\multicolumn{6}{|l|}{\hs{Gói 4. Quản lý thực phẩm trong tủ lạnh}}\\ \hline
UC24 & pantry & P01, P03 & \path{food_lots}; \path{user_settings} & — & —\\ \hline
UC25 & pantry \newline (+~catalog) & P02 (Q02, Q03, Y03, Y04) & \path{food_lots}, \path{stock_movements}; \path{foods}, \path{food_shelf_lives}; \path{media.objects} & \path{lot.changed} $\to$ đồng bộ; \path{food.changed} $\to$ ES (\path{index}) & \ref{fig:seq-them-thuc-pham}\\ \hline
UC26 & catalog & P08 & \path{foods}; Redis \path{barcode:{code}} & gọi Open Food Facts có ngắt mạch & \ref{fig:seq-ma-vach}\\ \hline
UC27 & pantry & P04 & \path{food_lots}, \path{stock_movements} & \path{lot.changed} $\to$ đồng bộ & —\\ \hline
UC28 & pantry & P05 & \path{food_lots}, \path{stock_movements} & \path{lot.changed} $\to$ đồng bộ & \ref{fig:seq-xac-nhan-nau}\\ \hline
UC29 & pantry & P06 & \path{food_lots}, \path{stock_movements} & \path{lot.changed} $\to$ đồng bộ & —\\ \hline
UC30 & pantry & P07 (M05) & \path{food_lots}; \path{user_settings} & — & \ref{fig:seq-goi-y-mon}\\ \hline
UC31 & notifi\-cation & (A14, A17, A18) & \path{notifications}, \path{deliveries}; \path{food_lots}, \path{user_settings}, \path{devices} & tác vụ định kỳ mỗi phút $\to$ \path{UserReminder} $\to$ \path{PushNotification} (\path{notify}) & \ref{fig:seq-nhac-han-dung-be}\\ \hline
\multicolumn{6}{|l|}{\hs{Gói 5. Quản lý công thức nấu ăn}}\\ \hline
UC32 & recipe & R02, R09 (Y03, Y04) & \path{recipes}, \path{recipe_ingredients} & \path{recipe.changed} $\to$ ES (\path{index}) & —\\ \hline
UC33 & recipe \newline (+~mealplan) & R04, R05 & \path{recipes}, \path{recipe_ingredients}; \path{meal_entries} & \path{recipe.changed} $\to$ ES; đồng bộ & —\\ \hline
UC34 & recipe \newline (+~pantry) & R03 (R01) & \path{recipes}, \path{recipe_ingredients}; \path{food_lots} & — & —\\ \hline
UC35 & recipe & R06, R07 & \path{favorite_recipes} & xóa \path{suggest:{sid}} & —\\ \hline
UC36 & recipe \newline (+~pantry, shopping) & R08 & \path{recipe_ingredients}; \path{food_lots}, \path{shopping_items} & đồng bộ & —\\ \hline
\multicolumn{6}{|l|}{\hs{Gói 6. Quản lý kế hoạch bữa ăn}}\\ \hline
UC37 & mealplan \newline (+~recipe) & M02 (R01) & \path{meal_entries}; \path{recipes} & đồng bộ & \ref{fig:seq-lap-lich}\\ \hline
UC38 & mealplan & M01 & \path{meal_entries} & — & —\\ \hline
UC39 & mealplan & M03, M04 & \path{meal_entries} & đồng bộ & —\\ \hline
UC40 & mealplan \newline (+~recipe, pantry) & M05 & \path{recipes}, \path{favorite_recipes}, \path{food_lots}; Redis \path{suggest:{sid}}, \path{cfg:all} & bộ đệm xóa khi có \path{lot.changed}, \path{recipe.changed}, \path{config.changed} & \ref{fig:seq-goi-y-mon}\\ \hline
UC41 & mealplan \newline (+~pantry) & M06 & \path{meal_entries}; \path{food_lots}, \path{stock_movements} & \path{meal.cooked}, \path{lot.changed} $\to$ đồng bộ & \ref{fig:seq-xac-nhan-nau}\\ \hline
\multicolumn{6}{|l|}{\hs{Gói 7. Tìm kiếm và phân loại thực phẩm}}\\ \hline
UC42 & search \newline (+~pantry) & Q01, Q02 & ES \path{foods}, \path{recipes}; \path{food_lots} (\path{StockQuery.SearchLots}) & chỉ mục nhận \path{food.changed}, \path{recipe.changed} (\path{index}) & \ref{fig:seq-tim-kiem-chi-muc}\\ \hline
UC43 & search \newline (+~pantry) & Q01, P01 (Q03) & ES \path{foods}; \path{food_categories}, \path{food_lots} & — & \ref{fig:seq-tim-kiem-chi-muc}\\ \hline
\multicolumn{6}{|l|}{\hs{Gói 8. Thống kê báo cáo}}\\ \hline
UC44 & report & B01 & \path{daily_space_stats}; \path{shopping_items}, \path{foods}, \path{food_categories} & làm mới hằng đêm & —\\ \hline
UC45 & report & B02 & \path{daily_space_stats}; \path{stock_movements}, \path{meal_entries} & làm mới hằng đêm & —\\ \hline
UC46 & report & B03 & \path{daily_space_stats}; \path{stock_movements}, \path{food_lots} & làm mới hằng đêm & \ref{fig:seq-bao-cao-xuat}\\ \hline
UC47 & report & B04, B05 & \path{report.exports}; MinIO \path{exports} & \path{ExportReport} (\path{report}) $\to$ \path{export.completed} $\to$ FCM (\path{notify}) & \ref{fig:seq-bao-cao-xuat}\\ \hline
\multicolumn{6}{|l|}{\hs{Gói 9. Quản trị hệ thống}}\\ \hline
UC48 & admin \newline (+~identity) & D01, D02 & \path{audit_logs}; \path{users}, \path{refresh_sessions}; Redis \path{auth:ver} & PUBLISH \path{user:{uid}}; thư thông báo (\path{notify}) & \ref{fig:trang-thai-tai-khoan}, \ref{fig:seq-quan-tri-cau-hinh}\\ \hline
UC49 & catalog & D03, D04, D11 & \path{foods}, \path{food_shelf_lives}; \path{audit_logs} & \path{food.changed} $\to$ ES (\path{index}); \path{food.merged} $\to$ đổi \path{food_id} ở pantry, shopping, recipe & \ref{fig:seq-duyet-thuc-pham}\\ \hline
UC50 & catalog & D05 & \path{food_categories}; \path{audit_logs} & đổi ETag danh mục & —\\ \hline
UC51 & catalog & D06 & \path{units}; \path{audit_logs} & đổi ETag danh mục & —\\ \hline
UC52 & recipe & D07 & \path{recipes}, \path{recipe_ingredients}; \path{audit_logs} & \path{recipe.changed} $\to$ ES (\path{index}) & —\\ \hline
UC53 & admin & D08, D09 (D10) & \path{system_configs}, \path{audit_logs}; Redis \path{cfg:all} & \path{config.changed} $\to$ Pub/Sub \path{cfg}, xóa \path{suggest:*} & \ref{fig:seq-quan-tri-cau-hinh}\\ \hline
\end{longtable}
}

\InputIfFileExists{chapters/c3/ht-dang-ky.tex}{}{}

% =====================================================================
\subsection{Tham gia nhóm gia đình (UC11)}\label{sssec:ht-tham-gia-nhom}

\uc{UC11} do dịch vụ \texttt{household.Service.JoinSpace} hiện thực qua G07 \texttt{POST} \path{/spaces/join}. Luồng gồm hai lời gọi, ứng với phần kiểm tra và xem trước (bước 4--6) và phần gia nhập sau khi người dùng xác nhận (bước 7--8) trong Bảng~\ref{tab:UC11}; Hình~\ref{fig:seq-tham-gia-nhom} thể hiện cộng tác của cả hai. Mã mời đến từ ô nhập, mã QR hoặc liên kết mời \path{https://dicho.example/join/{code}}; liên kết này được khai báo là Universal Links trên iOS và App Links trên Android nên hệ điều hành mở thẳng ứng dụng với mã điền sẵn (bước 3b).

\diagram[\textwidth]{c3-29-seq-tham-gia-nhom}{Biểu đồ tuần tự hiện thực use case Tham gia nhóm gia đình}{fig:seq-tham-gia-nhom}

Lời gọi thứ nhất mang trường \texttt{dryRun}: dịch vụ chỉ đọc, tra mã bằng một phép nối \texttt{invite\_codes} với \texttt{spaces} theo \texttt{code}, kiểm tra theo đúng thứ tự của Hình~\ref{fig:act-tham-gia-nhom} là \br{BR06}, \br{BR04} rồi \br{BR05} và trả tên nhóm, trưởng nhóm, số thành viên. Lời gọi xác nhận kiểm tra lại toàn bộ trong giao dịch ghi, vì giữa hai lần gọi mã có thể bị thu hồi hoặc nhóm có thể đã đủ người. Tình huống khó nhất là hai người nhập mã cùng lúc khi nhóm còn một chỗ: nếu mỗi giao dịch chỉ đếm rồi chèn, cả hai cùng thấy 9 người và nhóm thành 11 thành viên. Mức cô lập \texttt{READ COMMITTED} không ngăn được hiện tượng lệch ghi (write skew) này~\cite{kleppmann}. Sau khóa lũy đẳng, dịch vụ tra mã để biết \texttt{space\_id} rồi lấy \path{pg_advisory_xact_lock(hashtext(space_id::text))}, một khóa tư vấn gắn với giao dịch (bước 12 của hình). Giao dịch đến sau phải chờ (tối đa bằng \texttt{lock\_timeout}), rồi đọc lại mã và số thành viên bằng câu lệnh mới, thấy người vừa gia nhập và nhận \msg{MSG09}. G05--G10 lấy cùng khóa này nên mọi thay đổi thành viên, vai trò và mã mời của một nhóm được tuần tự hóa. Nhóm không khóa dòng \texttt{household.spaces} bằng \texttt{FOR UPDATE}, vì dòng này giữ bộ đếm \texttt{change\_seq} mà mọi giao dịch ghi của không gian đều cập nhật ở bước cuối; khóa nó từ đầu sẽ xếp hàng cả việc thêm mặt hàng hay nhập lô phía sau việc gia nhập và phá thứ tự khóa toàn cục. Hai không gian trùng giá trị băm chỉ phải chờ nhau một cách không cần thiết, không làm sai kết quả. Khóa tư vấn cũng không chặn được việc một người dùng hai thiết bị vào hai nhóm khác nhau, vì hai nhóm có hai khóa khác nhau; hàng rào cho \br{BR04} là chỉ mục \texttt{uq\_member\_family}, tức \texttt{UNIQUE(user\_id) WHERE space\_kind = 'family'}, với lỗi vi phạm 23505 được ánh xạ thành 422 mã 4004 kèm \msg{MSG08}. Mã mời chỉ hợp lệ khi \texttt{status = 'active'} và \texttt{expires\_at > now()}, không phụ thuộc tác vụ đánh dấu hết hạn hằng đêm. Giới hạn 10 lần mỗi giờ cho mỗi người (\texttt{rl:join:\{uid\}}) áp dụng cả cho lời gọi chạy thử, vì đó cũng là cách dò mã; với $32^8 \approx 1{,}1 \cdot 10^{12}$ tổ hợp, việc dò là không khả thi.

Cùng giao dịch còn cập nhật \texttt{last\_space\_id} qua cổng của identity để chuyển không gian làm việc (bước 8 của use case), chèn tác vụ \texttt{PushNotification} cho sự kiện \texttt{member.joined} và kết thúc bằng \texttt{space\_changes}. Sau commit, dịch vụ xóa khóa đệm \texttt{member:\{sid\}:\{uid\}} và phát \texttt{space:\{sid\}} để danh sách thành viên trên thiết bị khác cập nhật trong 3 giây (\nfr{NFR10}); worker báo cho trưởng nhóm theo cơ chế outbox của Hình~\ref{fig:seq-outbox-mua-hang} nếu người này bật thông báo hoạt động nhóm (\uc{UC08}). Các nhánh lỗi trả 422 mã 4004 kèm \br{BR06} (\msg{MSG07}), \br{BR04} (\msg{MSG08}) hoặc \br{BR05} (\msg{MSG09}); vượt giới hạn tần suất trả 429 mã 4290 kèm \msg{MSG36}. Nhờ khóa lũy đẳng, nếu phản hồi thành công bị mất, lần gửi lại nhận đúng phản hồi 200 đã lưu thay vì lỗi “đã là thành viên”; nếu lần gửi lại đến khi lần đầu chưa xong, ứng dụng nhận 409 mã 4002 và thử lại sau.

\InputIfFileExists{chapters/c3/ht-roi-nhom.tex}{}{}

% =====================================================================
\subsection{Phân công người mua (UC19)}\label{sssec:ht-phan-cong}

Phương thức \texttt{ShoppingService.AssignItem} hiện thực \uc{UC19} qua S10 \texttt{PUT} \path{.../shopping-items/{itemId}/assignee} để giao và S11 \texttt{DELETE} cùng đường dẫn để bỏ phân công (bước 3a), khi đó người được giao là rỗng. Hình~\ref{fig:seq-phan-cong} thể hiện vòng lặp phân công nhiều mặt hàng và nhánh bỏ phân công.

\diagram[\textwidth]{c3-30-seq-phan-cong}{Biểu đồ tuần tự hiện thực use case Phân công người mua}{fig:seq-phan-cong}

Mỗi lời gọi áp dụng cho một mặt hàng với phiên bản riêng trong \texttt{If-Match}, nhất quán với cơ chế đồng bộ theo thực thể (Hình~\ref{fig:seq-dong-bo-ngoai-tuyen}); chọn nhiều mặt hàng (bước 1) tương ứng một chuỗi lời gọi có kết quả độc lập, đúng yêu cầu ở bước 4b. Số việc của từng thành viên (bước 2) tính trên bản sao cục bộ. Quyền trưởng nhóm được kiểm tra hai lớp. Chính sách ở tầng ứng dụng đọc vai trò từ khóa đệm \texttt{member:\{sid\}:\{uid\}} để từ chối nhanh. Bộ đệm có thể chưa kịp phản ánh nhánh 4a, khi người dùng vừa chuyển quyền trên thiết bị khác, nên sau khóa lũy đẳng dịch vụ đọc lại dòng \texttt{space\_members} của người gọi và người được giao với \texttt{FOR SHARE}, theo thứ tự \texttt{user\_id} tăng dần. S10 không đổi thành viên nên không cần khóa tư vấn; G09 cập nhật chính các dòng này với \texttt{FOR UPDATE}, nên phân công và chuyển quyền luôn diễn ra theo một thứ tự xác định. Khóa trên dòng người được giao còn bảo vệ \br{BR07}: giao dịch rời nhóm hoặc xóa thành viên xóa dòng thành viên trước rồi mới gọi \texttt{MemberLeaveHook} để trả mặt hàng về Cần mua, nên phải chờ phân công kết thúc và thấy phân công đó; ngược lại, nếu người được giao đã rời nhóm, phép đọc không thấy dòng nào và lời gọi nhận 422 mã 4004 kèm \br{BR10}, \msg{MSG35}. Mặt hàng và danh sách bị khóa \texttt{FOR UPDATE} rồi mới so phiên bản; phân công là quyết định điều phối như đánh dấu đã mua nên không áp dụng quy tắc lệnh sau thắng, phiên bản lệch trả 409 mã 4002 kèm bản hiện tại và \msg{MSG35}. Điều kiện trạng thái là \texttt{to\_buy} hoặc \texttt{assigned} (\br{BR10}, Hình~\ref{fig:trang-thai-mat-hang-be}) và danh sách chưa Hoàn tất (\br{BR08}). Mỗi giao dịch kết thúc bằng \texttt{space\_changes}, sau commit phát \texttt{space:\{sid\}}; S11 đi qua cùng các bước kiểm tra vai trò và phiên bản.

Thông báo phải gộp: khi trưởng nhóm giao năm mặt hàng liên tiếp, người được giao chỉ nhận một thông báo “Bạn được giao mua $k$ mặt hàng trong danh sách <tên>” (bước 6). Đồng thời, module notification không được đọc bảng của shopping (mục~\ref{sec:kien-truc}), nên $k$ phải nằm sẵn trong nội dung sự kiện. Hai yêu cầu được dung hòa bằng một tác vụ gom của chính shopping. Giao dịch S10 chèn tác vụ \texttt{AssignmentDigest} duy nhất theo cặp người được giao và danh sách, lịch chạy lùi 5 giây; bốn lần chèn sau trùng tác vụ đang chờ nên bị River bỏ qua. Khi chạy, tác vụ đếm các mặt hàng của danh sách được giao cho người đó từ mốc đợt giao, rồi chèn \texttt{PushNotification} mang sự kiện \texttt{item.assigned} với $k$, tên danh sách và mốc đợt. Worker thông báo tôn trọng thiết lập \texttt{assignment\_notify\_on} (\uc{UC08}), lưu thông báo vào hộp thư với \texttt{dedupe\_key} gồm người nhận, danh sách và mốc đợt để lần thử lại không sinh bản trùng, rồi đẩy FCM tới từng thiết bị; lỗi tạm thời được thử lại tối đa ba lần trên hàng đợi \texttt{notify} (bước 6a, \nfr{NFR12}). Tự giao cho chính mình thì không đẩy; bỏ phân công dùng cùng cơ chế với sự kiện \texttt{item.unassigned} để báo người được giao trước đó. Về lỗi, người không phải trưởng nhóm nhận 403 mã 3007 (\msg{MSG13}) ở cả hai lớp kiểm tra; mặt hàng Đã mua, Đã bỏ qua, Đã nhập tủ hoặc thuộc danh sách Hoàn tất nhận 422 mã 4004 kèm \br{BR10} hoặc \br{BR08}. Ứng dụng tổng hợp các phản hồi, hiển thị \msg{MSG10} cùng số mặt hàng bị bỏ qua.

% =====================================================================
\subsection{Sinh danh sách mua sắm từ kế hoạch bữa ăn (UC22)}\label{sssec:ht-sinh-ds}

Theo Bảng~\ref{tab:UC22}, \uc{UC22} có hai pha: tính đề xuất (bước 2--5) và tạo danh sách sau khi người dùng chỉnh sửa (bước 6--7). Cả hai pha đi qua S15 \texttt{POST} \path{.../shopping-lists/generate}, do \texttt{MealHandler} chuyển cho \texttt{MealPlanService.GenerateList} của module mealplan; lời gọi thứ nhất mang \texttt{dryRun}, lời gọi thứ hai mang các dòng đã chọn, như Hình~\ref{fig:seq-sinh-ds-ke-hoach}. Ba cổng của module khác được vẽ gộp thành một bên tham gia.

\diagram[\textwidth]{c3-31-seq-sinh-ds-ke-hoach}{Biểu đồ tuần tự hiện thực use case Sinh danh sách mua sắm từ kế hoạch bữa ăn}{fig:seq-sinh-ds-ke-hoach}

Dù đường dẫn nằm dưới tài nguyên danh sách, dịch vụ thuộc mealplan vì phép tính cần dữ liệu của bốn module. Mealplan được phép gọi cổng \texttt{RecipeQuery} của recipe, \texttt{StockQuery} của pantry và \texttt{ListComposer} của shopping (hai phương thức \texttt{PendingQuantities} và \texttt{Compose}), còn recipe đã phụ thuộc shopping vì \uc{UC36}; đặt \uc{UC22} trong shopping sẽ tạo chu trình shopping $\to$ recipe $\to$ shopping. Khoảng ngày phải thỏa \texttt{from} $\le$ \texttt{to} và nằm trong cửa sổ từ \texttt{\$today} đến số ngày của tham số \texttt{mealplan.horizonDays} (mặc định 28) theo \br{BR18}, nếu không trả 400 mã 2002 (\msg{MSG02}). Pha xem trước không ghi gì: bốn phép đọc gồm món chưa nấu có công thức, nguyên liệu nhân theo tỷ lệ khẩu phần, lượng đang chờ mua và lô còn hạn chạy trong một giao dịch \texttt{REPEATABLE READ READ ONLY}. Nhờ ảnh chụp chung, khi thành viên khác đang xác nhận nấu (\uc{UC41}) hay nhập tủ (\uc{UC21}), phép tính không thấy trạng thái nửa trước nửa sau như lô đã bị trừ mà món vẫn Dự kiến; giao dịch chỉ đọc cũng không giữ khóa dòng nên không cản ai. Lượng đang chờ mua gồm mặt hàng Cần mua, Đã phân công và Đã mua nhưng chưa nhập tủ, thuộc danh sách chưa Hoàn tất có ngày mua không sau ngày nấu (\br{BR20}). Thành phần miền \texttt{ShortageCalculator.Compute} là hàm thuần, kiểm thử đơn vị được trực tiếp bằng ví dụ cà rốt của Chương~\ref{chap:dac-ta}; thuật toán đã trình bày ở mục~\ref{sssec:kt-goi-y}. Món tên tự do được trả trong danh sách riêng để ghi chú (bước 2a). Khi không thiếu nguyên liệu nào, S15 trả 200 mã 1000 kèm \msg{MSG15} và danh sách rỗng (bước 5a), vì đây là thông tin chứ không phải lỗi.

Lời gọi xác nhận không tính lại phần thiếu, vì đề xuất chỉ mang tính tư vấn và người dùng được chỉnh ở bước 6. Trước khi mở giao dịch, dịch vụ kiểm tra số lượng (\br{BR14}, \msg{MSG16}) và ngày dự kiến mua không trước \texttt{\$today} (\br{BR08}, \msg{MSG02}); khi lệnh được phát lại từ hàng đợi ngoại tuyến, mốc so sánh là ngày của header \texttt{X-Client-Time} nếu thời điểm đó trong vòng 7 ngày. Giao dịch ghi bắt đầu bằng khóa lũy đẳng rồi giao việc tạo cho cổng \texttt{ListComposer.Compose}, nhờ đó quy tắc gộp chỉ cài đặt ở module shopping. Tạo mới (bước 7) chèn danh sách với \texttt{origin = meal\_plan}, trạng thái \texttt{drafting} cùng các mặt hàng \texttt{to\_buy}. Với nhánh gộp (bước 6a), ứng dụng lấy các danh sách ứng viên qua S01 lọc theo trạng thái và khoảng ngày; \texttt{Compose} khóa \texttt{FOR UPDATE} dòng danh sách đích để tuần tự hóa với thành viên đang thêm mặt hàng (\uc{UC17}), danh sách đã Hoàn tất làm giao dịch bị hủy với 422 mã 4004 kèm \br{BR08}, còn lại gộp theo \br{BR09} và trả \msg{MSG11}. Giao dịch kết thúc bằng \texttt{space\_changes}, sau commit phát \texttt{space:\{sid\}}; khóa lũy đẳng bảo đảm gửi lại không tạo hai danh sách.

% =====================================================================
\subsection{Thêm thực phẩm vào tủ lạnh (UC25)}\label{sssec:ht-them-thuc-pham}

Phương thức \texttt{PantryService.AddLot} hiện thực \uc{UC25} qua P02 \texttt{POST} \path{/spaces/{sid}/lots}, với các API hỗ trợ trong cùng luồng: P08 cho nhánh quét mã vạch (bước 2a), Q02 \texttt{GET} \path{/spaces/{sid}/foods/suggest} cho gợi ý tên (bước 3--4), Q03 cho danh mục loại, đơn vị và vị trí lưu trữ, Y03 và Y04 cho ảnh của thực phẩm mới (bước 3a). Hình~\ref{fig:seq-them-thuc-pham} thể hiện luồng đầy đủ; nhánh mã vạch được tham chiếu sang Hình~\ref{fig:seq-ma-vach}.

\diagram[\textwidth]{c3-32-seq-them-thuc-pham}{Biểu đồ tuần tự hiện thực use case Thêm thực phẩm vào tủ lạnh}{fig:seq-them-thuc-pham}

Q02 tìm trên Elasticsearch với bộ lọc phạm vi bắt buộc, chỉ gồm thực phẩm dùng chung và thực phẩm riêng của không gian hiện hành (\nfr{NFR09}, \br{BR23}) và trả kèm đơn vị, vị trí mặc định theo loại cùng số ngày hạn dùng theo từng vị trí lấy từ \texttt{catalog.food\_shelf\_lives}, nên ứng dụng điền sẵn ngày hết hạn ở bước 4 mà không cần thêm lượt gọi. P02 dùng được khi ngoại tuyến (\nfr{NFR11}) nhờ id UUIDv7 do thiết bị sinh và \texttt{Idempotency-Key}; khi phát lại, ràng buộc ngày mua không sau hôm nay được xét theo ngày của \texttt{X-Client-Time}. Dịch vụ kiểm tra dữ liệu vào trước mọi truy vấn, đúng bước 6: trường bắt buộc, số lượng (\br{BR14}), ngày mua không sau \texttt{\$today} và ngày hết hạn không trước ngày mua (\br{BR11}); ràng buộc cuối được lặp lại bằng \texttt{CHECK} \texttt{ck\_lot\_dates} trên \texttt{food\_lots}. Trong giao dịch, sau khóa lũy đẳng, thực phẩm phải là thực phẩm dùng chung đang hoạt động hoặc thực phẩm riêng của chính không gian; thực phẩm của không gian khác làm giao dịch bị hủy với 404 như không tồn tại (\br{BR23}). Ngày hết hạn bỏ trống được đề xuất theo \br{BR11} bằng ngày mua cộng số ngày của cặp thực phẩm, vị trí; thực phẩm chưa có hạn dùng gợi ý thì trường này bắt buộc. Ngày hết hạn đã qua (bước 7a) không bị chặn: ứng dụng hỏi xác nhận, máy chủ vẫn lưu và trả trạng thái Đã hết hạn được suy diễn lúc đọc (\br{BR12}).

Thực phẩm mới (bước 3a) được gửi trong trường \texttt{newFood} và tạo trong cùng giao dịch qua cổng \texttt{catalog/api.CreateSpaceFood}: bản ghi \texttt{catalog.foods} có \texttt{scope = space}, chủ sở hữu là không gian hiện hành và \texttt{status = pending\_review}; chỉ mục duy nhất từng phần theo không gian, loại và tên ngăn hai thành viên tạo trùng, khi xung đột thì dùng lại bản ghi đã có. Sự kiện \texttt{food.changed} sinh tác vụ \texttt{IndexFood} để thực phẩm xuất hiện trong gợi ý của riêng không gian này và trong hàng chờ duyệt của \uc{UC49} (Hình~\ref{fig:seq-duyet-thuc-pham}). Ảnh không đi qua API: Y03 nhận \texttt{\{purpose: food, targetId\}}, kiểm tra quyền trên đối tượng đích, tạo \texttt{media.objects} ở trạng thái \texttt{pending} và trả URL \texttt{PUT} ký sẵn (5 phút, tối đa 5 MB, JPEG hoặc PNG). Ứng dụng tải thẳng lên MinIO rồi gọi Y04; API dùng lệnh \texttt{HEAD} để xác nhận đối tượng có thật, đúng kích thước và kiểu, rồi chèn tác vụ xử lý ảnh cho sự kiện \texttt{media.uploaded} để \texttt{MediaWorker} nén ảnh và tạo ảnh thu nhỏ WebP. Giao dịch ghi chèn \texttt{food\_lots} với \texttt{unit\_factor} là ảnh chụp hệ số quy đổi, để việc sửa đơn vị ở \uc{UC51} không làm đổi lượng của lô cũ, cùng biến động \texttt{stock\_in} làm đầu vào tỷ lệ lãng phí của \uc{UC46}; lô nhập tay không có \texttt{total\_cost} vì chỉ S14 biết đơn giá. Giao dịch kết thúc bằng \texttt{space\_changes}; sau commit, dịch vụ xóa bộ đệm gợi ý \texttt{suggest:\{sid\}} và phát \texttt{space:\{sid\}}; phản hồi 201 trả lô kèm trạng thái hạn dùng theo ngưỡng $N$ của người gọi. Thiếu trường trả 400 mã 2001 (\msg{MSG01}); số lượng sai và ngày hết hạn trước ngày mua trả 400 mã 2002 (\msg{MSG16}, \msg{MSG14}). Lỗi tra mã vạch (mã 4001 hoặc 5004, \msg{MSG23}) đưa người dùng về biểu mẫu để nhập tay, đúng bước 4b của \uc{UC26}.

% =====================================================================
\subsection{Gợi ý món ăn (UC40)}\label{sssec:ht-goi-y-mon}

Phương thức chỉ đọc \texttt{MealPlanService.Suggest} hiện thực \uc{UC40} qua M05 \texttt{GET} \path{/spaces/{sid}/recipe-suggestions} với các tham số \texttt{date}, \texttt{mealSlot} (\texttt{breakfast}, \texttt{lunch} hoặc \texttt{dinner}), số khẩu phần và \texttt{expiringOnly}. Cùng một API phục vụ ba điểm vào: thẻ “Gợi ý cho bạn” của \uc{UC37} (bước 2a), màn hình gợi ý độc lập và bước 3 của \uc{UC30}. Ở \uc{UC30}, ứng dụng ghép P07 với M05 có \texttt{expiringOnly = true} vì pantry không được gọi mealplan. Hình~\ref{fig:seq-goi-y-mon} thể hiện luồng xử lý.

\diagram[\textwidth]{c3-33-seq-goi-y-mon}{Biểu đồ tuần tự hiện thực use case Gợi ý món ăn từ thực phẩm sẵn có}{fig:seq-goi-y-mon}

Theo điều kiện sau của Bảng~\ref{tab:UC40}, kết quả gợi ý được lưu 10 phút hoặc đến khi dữ liệu đầu vào đổi, trong Redis hash \texttt{suggest:\{sid\}} với trường gồm người dùng, ngày, bữa, số khẩu phần và cờ \texttt{expiringOnly}, vì tồn kho thuộc không gian còn ngưỡng $N$ và hệ số $F$ thuộc từng người. Hash bị xóa sau commit của mọi giao dịch phát \texttt{lot.changed} hay \texttt{recipe.changed}, khi quản trị viên đổi tham số gợi ý (\texttt{config.changed}) và khi A13 đổi $N$ của một người (xóa ở mọi không gian của người đó); thời hạn sống chặn trên độ cũ nếu lệnh xóa thất bại. Khi bộ đệm trống, dịch vụ đọc $\theta$, trọng số và số món tối đa từ \texttt{cfg:all}, nên thay đổi ở \uc{UC53} có hiệu lực ngay lần tính sau. Ba truy vấn trong một giao dịch chỉ đọc lấy ngưỡng $N$ của người xem, tồn kho khả dụng qua \texttt{StockQuery.Available} (lô \texttt{active} có \texttt{expires\_on} $\ge$ \texttt{\$today}) và qua \texttt{RecipeQuery} các công thức ở trạng thái \texttt{published} có ít nhất một nguyên liệu trong tủ, lọc theo bữa, cùng danh sách yêu thích; công thức mẫu đang ẩn hay bản nháp không bao giờ được gợi ý. Thành phần miền \texttt{SuggestionEngine.Rank} chấm điểm theo \br{BR19} (mục~\ref{sssec:kt-goi-y}) với mục tiêu 2 giây (\nfr{NFR09}); với \texttt{expiringOnly}, chỉ giữ các món có $E > 0$, tức dùng ít nhất một lô sắp hết hạn, và ứng dụng hiển thị ba món đầu. Tủ lạnh trống (bước 2a) trả công thức mẫu được yêu thích nhiều nhất kèm lời nhắc thêm thực phẩm; không món nào đạt ngưỡng (bước 4a) trả ba món có $R$ cao nhất kèm \msg{MSG31}, cả hai đều mang mã 1000. Tham số sai trả 400 mã 2002 (\msg{MSG02}). Là phép đọc thuần, M05 không dùng khóa lũy đẳng và không giữ khóa dòng.

% =====================================================================
\subsection{Lên lịch bữa ăn (UC37)}\label{sssec:ht-lap-lich}

Sau khi người dùng chọn món ở một trong ba thẻ, gồm Gợi ý (M05, Hình~\ref{fig:seq-goi-y-mon}), Yêu thích và Tất cả công thức (R01, chỉ trả công thức \texttt{published}), M02 \texttt{POST} \path{/spaces/{sid}/meals} gọi \texttt{MealPlanService.PlanMeals} để ghi các món vào lịch, như Hình~\ref{fig:seq-lap-lich}.

\diagram[\textwidth]{c3-40-seq-lap-lich}{Biểu đồ tuần tự hiện thực use case Lên lịch bữa ăn}{fig:seq-lap-lich}

M02 nhận một ngày, một bữa và mảng món; mỗi món có id do thiết bị sinh, tham chiếu công thức hoặc tên tự do (bước 3a) và số khẩu phần, nên chọn nhiều món ở bước 3 được ghi trong một giao dịch. Trước mọi truy vấn, dịch vụ kiểm tra \br{BR18}: ngày nằm từ \texttt{\$today} đến hết cửa sổ lập kế hoạch đọc từ tham số \texttt{mealplan.horizonDays} (mặc định 28 ngày, miền 7–60, nếu vi phạm trả \msg{MSG28}), và số khẩu phần từ 1 đến 20. Khi lệnh được phát lại từ hàng đợi ngoại tuyến, \br{BR18} được xét theo ngày của \texttt{X-Client-Time} nếu thời điểm đó trong vòng 7 ngày và không vượt giờ máy chủ; ngoài khoảng đó máy chủ trả 422 mã 4004 và ứng dụng giữ bản nháp để người dùng sửa. Trong giao dịch, sau khóa lũy đẳng, \texttt{RecipeQuery} xác nhận công thức là công thức mẫu đang công bố hoặc thuộc chính không gian; công thức ngoài phạm vi làm giao dịch bị hủy với 404 mã 4001. Món được chèn ở trạng thái \texttt{planned}, giao dịch kết thúc bằng \texttt{space\_changes} và sau commit phát \texttt{space:\{sid\}}. Món trùng trong cùng bữa (bước 5a) được ứng dụng phát hiện trên bản sao cục bộ; nếu người dùng chọn tăng khẩu phần, ứng dụng gọi M03 thay cho M02. Sau commit, dịch vụ đối chiếu nguyên liệu bằng \texttt{RecipeQuery.Ingredients} và \texttt{StockQuery.Available}, cũng là thành phần dùng chung với \uc{UC34}, và trả nhãn đủ hoặc thiếu $k$ nguyên liệu cho từng món (bước 5); món tên tự do không tham gia đối chiếu. Phép đối chiếu nằm ngoài giao dịch ghi để giữ giao dịch ngắn. Ngày ngoài khoảng trả 400 mã 2002 với \msg{MSG28}, khẩu phần sai trả \msg{MSG02}; M03 áp dụng cùng quy tắc khi di chuyển món (\uc{UC39}), còn món đã nấu không được sửa.

% =====================================================================
\subsection{Báo cáo lãng phí và xuất báo cáo (UC46, UC47)}\label{sssec:ht-bao-cao}

\texttt{ReportService} hiện thực \uc{UC46} qua B03 \texttt{GET} \path{.../reports/waste} và \uc{UC47} qua cặp B04, B05: B04 nhận yêu cầu xuất và trả ngay, tác vụ River \texttt{ExportReport} kết xuất tệp vào MinIO, B05 trả trạng thái và cấp URL tải ký sẵn. Report chỉ đọc dữ liệu của module khác và không gọi cổng nào, nên Hình~\ref{fig:seq-bao-cao-xuat} vẽ handler và dịch vụ thành một bên tham gia.

\diagram[\textwidth]{c3-34-seq-bao-cao-xuat}{Biểu đồ tuần tự hiện thực use case Xem báo cáo lãng phí và Xuất báo cáo}{fig:seq-bao-cao-xuat}

Bước 4 của \uc{UC46} kiểm tra \br{BR21}; khoảng sai hoặc dài hơn 12 tháng trả 400 mã 2002 (\msg{MSG25}). Report là module mô hình đọc (read model): ngoại lệ có chủ đích của quy tắc phụ thuộc cho phép nó \texttt{SELECT} qua các khung nhìn \texttt{report.v\_*} đặt trên bảng của shopping, pantry, mealplan và catalog (mục~\ref{sec:kien-truc}). Nguồn số liệu được thiết kế lai để chi phí không tăng theo độ dài khoảng. Phần lịch sử đến hết hôm qua đọc từ \texttt{report.daily\_space\_stats}, khung nhìn cụ thể hóa làm mới hằng đêm, mỗi dòng ứng với một không gian, một ngày, một loại thực phẩm và một nhóm đại lượng, nên khoảng 12 tháng chỉ phải cộng vài nghìn dòng đã tổng hợp thay vì quét toàn bộ biến động kho. Các chỉ số không tổng hợp sẵn gồm số lô bị loại bỏ trên số lô nhập tủ (tỷ lệ lãng phí theo Bảng~\ref{tab:chi-so-bao-cao}), phân bổ theo lý do, năm thực phẩm bị bỏ nhiều nhất và số liệu của \texttt{\$today}; chúng được truy vấn trực tiếp từ \texttt{pantry.stock\_movements} qua chỉ mục (\texttt{space\_id}, \texttt{occurred\_at}). Giá trị lãng phí của một lần loại bỏ bằng \texttt{total\_cost} $\times$ lượng loại bỏ / \texttt{initial\_qty}, trong đó \texttt{total\_cost} của lô là số lượng thực tế nhân đơn giá ghi nhận khi nhập tủ từ mặt hàng nguồn (S14). Cách tính theo tỷ lệ cho đúng đơn vị kể cả khi S14 đổi đơn vị so với mặt hàng; lô không có \texttt{total\_cost} chỉ tính theo số lượng và được đếm riêng (\br{BR21}). Hai phần đọc chung một giao dịch chỉ đọc nên các thẻ chỉ số nhất quán; truy vấn cùng độ dài cho kỳ trước cho mức thay đổi ở bước 6, và cờ \texttt{empty} phục vụ bước 5a. Người không phải thành viên, kể cả quản trị viên, nhận 404 mã 4001 (\br{BR23}).

\uc{UC47} được xử lý bất đồng bộ, vì kết xuất PDF có biểu đồ tốn vài giây CPU, nếu giữ trong yêu cầu sẽ chiếm tiến trình API, vỡ mục tiêu 95\% API dưới 1 giây (\nfr{NFR09}) và dễ đứt trên mạng di động. B04 nhận loại báo cáo, khoảng thời gian và định dạng PDF hoặc CSV (bước 1). Sau khóa lũy đẳng, giao dịch chèn một dòng vào bảng \texttt{report.exports} do module report sở hữu, gồm không gian, người yêu cầu, loại, khoảng, định dạng và trạng thái \texttt{pending}, cùng tác vụ \texttt{ExportReport} trên hàng đợi \texttt{report}, duy nhất theo người dùng, loại, khoảng và định dạng để chạm nhiều lần không sinh nhiều tệp; API trả ngay 202 mã 1001 kèm \texttt{exportId}. Bảng \texttt{media.objects} chỉ dành cho ảnh. Worker gọi lại đúng truy vấn của B01--B03 nên số liệu trong tệp trùng với màn hình, kết xuất PDF hoặc CSV mã hóa UTF-8 (bước 2). Khi thành công, tệp được ghi vào bucket \texttt{exports} có vòng đời 7 ngày; trong một giao dịch, dòng xuất chuyển sang \texttt{ready} với \texttt{object\_key}, \texttt{expires\_at}, và tác vụ \texttt{PushNotification} cho sự kiện \texttt{export.completed} được chèn để module notification báo “Báo cáo đã sẵn sàng”. Kết xuất lỗi (bước 2a) được River thử lại tối đa ba lần; sau lần cuối, dòng xuất chuyển sang \texttt{failed} kèm lỗi và người dùng cũng được báo qua \texttt{PushNotification}. Ứng dụng gọi B05 khi nhận thông báo hoặc thăm dò thưa dần khi màn hình còn mở. Sau bước kiểm tra thành viên chung, B05 lọc theo \texttt{id}, \texttt{space\_id} và \texttt{requested\_by}, nên \texttt{exportId} không tồn tại, thuộc không gian khác hoặc của người khác đều nhận 404 mã 4001. Phản hồi phân theo trạng thái: \texttt{pending} để ứng dụng tiếp tục chờ; \texttt{ready} kèm URL \texttt{GET} ký sẵn 15 phút, ứng dụng tải thẳng từ MinIO và mở bảng chia sẻ của hệ điều hành (bước 3); \texttt{failed} để ứng dụng hiện nút thử lại, tức gửi một yêu cầu B04 mới.

\InputIfFileExists{chapters/c3/ht-quan-tri.tex}{}{}

% =====================================================================
\subsection{Các use case còn lại}\label{sssec:ht-con-lai}

Đăng nhập và làm mới phiên (\uc{UC03}) có biểu đồ ở mục~\ref{ssec:xac-thuc-phien} (Hình~\ref{fig:seq-dang-nhap-lam-moi}). Các use case khác đã có biểu đồ ở mục~\ref{ssec:ky-thuat} gồm \uc{UC20} (Hình~\ref{fig:seq-outbox-mua-hang}), \uc{UC21} (Hình~\ref{fig:seq-nhap-tu}), \uc{UC26} (Hình~\ref{fig:seq-ma-vach}), \uc{UC31} (Hình~\ref{fig:seq-nhac-han-dung-be}), \uc{UC41} (Hình~\ref{fig:seq-xac-nhan-nau}), \uc{UC42} và \uc{UC43} (Hình~\ref{fig:seq-tim-kiem-chi-muc}); trong đó \uc{UC21} gọi cổng \texttt{StockIn.AddLots} và \uc{UC41} gọi \texttt{StockConsumer.ConsumeForMeal} của pantry trong cùng giao dịch với module gọi. \uc{UC28} dùng chung cơ chế FEFO của Hình~\ref{fig:seq-xac-nhan-nau} cho một thực phẩm: P05 gọi \texttt{PantryService.Consume}, khóa các lô \texttt{active} có \texttt{expires\_on} $\ge$ \texttt{\$today} theo thứ tự (\texttt{expires\_on}, \texttt{id}) rồi trừ dần theo \br{BR15}; khi chọn đích danh một lô (bước 3a) chỉ lô đó bị khóa, lô hết hạn chỉ bị trừ khi được chọn tường minh (bước 4b), và lượng vượt lượng còn trả 422 mã 2003 (\msg{MSG17}). P04 của \uc{UC27} không áp dụng lệnh sau thắng cho lượng còn lại: \texttt{If-Match} lệch trả 409 mã 4002, thay đổi lượng ghi biến động \texttt{adjust}, lượng lớn hơn lượng ban đầu trả 400 mã 2002 và lượng bằng 0 chuyển lô sang Đã dùng hết; vị trí, ngày hết hạn và ghi chú thì hợp nhất theo trường kèm \texttt{overwrittenFields}. P06 của \uc{UC29} bắt buộc lý do và ghi biến động \texttt{discard} cho \uc{UC46}.

Ở gói tài khoản và nhóm, A15 \texttt{DELETE} \path{/me/devices/{deviceId}} gỡ mã FCM khi đăng xuất (\uc{UC05}). A10 cập nhật hồ sơ (\uc{UC06}) ghi \texttt{space\_changes} với thực thể thành viên ở mọi không gian của người dùng để tên và ảnh đại diện mới đến được các thành viên khác. A11 đổi mật khẩu (\uc{UC07}) thu hồi mọi phiên khác, tăng \texttt{token\_version} rồi cấp cặp mã mới cho phiên hiện tại, nên thiết bị khác bị đăng xuất ngay còn người dùng không phải đăng nhập lại (bước 3). A12 trả thiết lập kèm khung giờ nhắc \texttt{reminder.runAt} và ngưỡng $N$ mặc định ở dạng chỉ đọc; A13 đổi $N$ thì xóa bộ đệm gợi ý của mọi không gian của người đó (\uc{UC08}). A16 xóa tài khoản theo \nfr{NFR18}: người dùng nhập lại mật khẩu, trưởng nhóm của nhóm còn thành viên khác bị từ chối với 422 mã 4004 kèm \msg{MSG24} (\br{BR07}); ngay lập tức mọi phiên bị thu hồi, \texttt{deleted\_at} được đặt, đăng nhập bị chặn và thiết bị bị xóa; sau 30 ngày, tác vụ nền ẩn danh email và họ tên, xóa ảnh, dữ liệu không gian cá nhân và tài liệu Elasticsearch, còn dữ liệu đóng góp trong nhóm được giữ với tên “Người dùng đã xóa”. G02 tạo nhóm (\uc{UC09}) dùng chung chỉ mục \br{BR04} với \uc{UC11}; trường \texttt{copyPersonalData} gồm danh sách, lô và công thức làm giao dịch chèn tác vụ \texttt{CopySpaceData} lũy đẳng theo không gian để sao chép sau commit, còn G04 đổi tên, ảnh nhóm. G05 tạo mã mời (\uc{UC10}) chỉ dành cho trưởng nhóm, lấy khóa tư vấn của không gian, từ chối không gian cá nhân với 422 mã 4004, thu hồi mã cũ và chèn mã mới hết hạn sau \texttt{invite.ttlHours} trong một giao dịch (\br{BR06}); Caddy phục vụ hai tệp \path{/.well-known/apple-app-site-association} và \path{/.well-known/assetlinks.json} để liên kết mời mở được ứng dụng (mục~\ref{sec:trien-khai}).

G03 (\uc{UC12}) trả thành viên, vai trò và \texttt{recentActivity} gồm 20 thay đổi gần nhất trong \texttt{sync.space\_changes} kèm tên người thực hiện (bước 3); \uc{UC13}--\uc{UC15} đã được trình bày ở Hình~\ref{fig:seq-roi-nhom-chuyen-quyen}.

Ở các gói còn lại, S01 (\uc{UC16}, \uc{UC23}) lọc theo trạng thái, khoảng ngày và gom theo tháng, phản hồi có tổng chi tiêu và người mua. S07 khóa danh sách trước khi gộp theo \br{BR09} (\msg{MSG11}); bỏ qua mặt hàng là lệnh S08 đặt \texttt{status = skipped}; mặt hàng Đã nhập tủ chỉ đọc (\msg{MSG32}). Danh sách Hoàn tất chặn sửa nội dung với 422 mã 4004 (\br{BR08}), trừ ba lệnh: S08 sửa số lượng thực tế và đơn giá của mặt hàng Đã mua, S13 hoàn tác đưa danh sách về Đang mua, và S14 nhập tủ. Công thức thuộc không gian: ở không gian cá nhân, đây là “bộ sưu tập cá nhân” của Chương~\ref{chap:dac-ta}. Chỉ người tạo hoặc trưởng nhóm được sửa, xóa công thức (\uc{UC33}); sửa công thức mẫu trả 403 mã 3007. R09 sao chép một công thức mẫu thành bản của không gian, ghi \texttt{origin\_recipe\_id} (\uc{UC32}, \br{BR17}). Trạng thái nháp, công bố và ẩn chỉ dùng cho công thức mẫu (\uc{UC52}); công thức của không gian luôn ở trạng thái công bố. R06 (\uc{UC35}) chỉ nhận công thức mẫu đã công bố hoặc công thức của không gian mà người gọi là thành viên, ngược lại trả 404 mã 4001. Xóa công thức chuyển các món Dự kiến thành món tên tự do qua \texttt{RecipeDeleteHook}. R03 (\uc{UC34}) và R08 (\uc{UC36}) đối chiếu nguyên liệu với tồn kho qua \texttt{StockQuery}; R08 giao phần thiếu cho \texttt{ListComposer}.

P01 và Q01 (\uc{UC43}) nhận bộ lọc loại, vị trí, trạng thái hạn dùng và từ khóa, phản hồi kèm số lượng theo loại (facets); Q01 tìm danh mục chung, thực phẩm riêng và công thức của không gian trên Elasticsearch với bộ lọc phạm vi bắt buộc, tìm lô trong không gian qua \texttt{StockQuery.SearchLots} trên PostgreSQL và chuyển hẳn sang PostgreSQL khi Elasticsearch lỗi. B01 (\uc{UC44}) gom theo ngày, tuần hoặc tháng; B01, B02 (\uc{UC45}) dùng cùng nguồn lai như B03. Ngoài \uc{UC48}, \uc{UC49} và \uc{UC53} đã trình bày ở trên, các use case quản trị theo cùng khuôn: mọi API yêu cầu \texttt{system\_role = admin} và ghi \texttt{admin.audit\_logs} gồm giá trị cũ và mới trong cùng giao dịch. D05 và D06 (\uc{UC50}, \uc{UC51}) xóa được bản ghi chưa được tham chiếu; bản ghi đang được tham chiếu trả 422 mã 4004 kèm \br{BR22} và \msg{MSG22} để ứng dụng đề nghị chuyển sang Ngừng sử dụng, mọi thay đổi đều làm đổi ETag của Q03. Đơn vị cơ sở là duy nhất trong mỗi nhóm đại lượng và nhóm đại lượng của đơn vị đang được tham chiếu không đổi được (\br{BR16}). D07 (\uc{UC52}) quản lý công thức mẫu; mỗi lần công bố hay ẩn phát \texttt{recipe.changed} để đánh lại chỉ mục.
